\section{Estado global, instrumentos locales y correlaciones triádicas de Bell}
\label{chap:bell-local-global}

La cuestión planteada por Bell \cite{Bell1964} relaciona la descripción local de los
resultados con la estructura global de sus correlaciones. Pregunta si una familia de
probabilidades conjuntas admite una descomposición en respuestas locales
condicionadas por un mismo estado previo. En HMT esta pregunta dispone de un dominio natural: el estado enriquecido conserva la
genealogía, la orientación, el acarreo y la memoria de frontera, mientras cada
instrumento publica una coordenada local. La riqueza del estado común y la forma
probabilística de sus lecturas constituyen datos matemáticos distintos. La primera
organiza la procedencia de la correlación; la segunda determina su posición respecto
del politopo local de Bell.

El propósito de este capítulo es recorrer esa cadena completa. Partiremos de una
preparación global sobre el espacio de estados TPK, definiremos los instrumentos de
las dos alas, recuperaremos la cota CHSH y demostraremos su estabilidad bajo el
enriquecimiento del estado compartido; pasaremos después al alfabeto ternario propio
del TRIT. En ese nuevo
dominio la unidad informativa es el trit y la desigualdad pertinente es CGLMP. La
estructura elíptica interna proporciona además una realización real de la transformada
de Fourier de orden tres mediante el operador \(J_E^2=-I\). El resultado es
una formulación precisa de la relación entre totalidad holográfica, lectura local y
correlación cuántica.

\subsection{Preparación enriquecida e instrumentos de frontera}

Sea \(X_{\rm enr}\) el espacio de estados enriquecidos de TPK y sea
\(\mu\) una medida de probabilidad que describe la preparación común. Un elemento
\(x\in X_{\rm enr}\) comprende el estado local junto con su historia de transporte.
La lectura se introduce mediante instrumentos. Para dos
alas, Alice y Bob, con conjuntos de ajustes \(\mathcal A\) y \(\mathcal B\),
escribimos
\[
 \mathcal I_A^a:X_{\rm enr}\longrightarrow\mathcal D(\mathcal O_A),
 \qquad
 \mathcal I_B^b:X_{\rm enr}\longrightarrow\mathcal D(\mathcal O_B),
\]
donde \(\mathcal D(\mathcal O)\) denota el simplex de distribuciones sobre el
alfabeto de resultados \(\mathcal O\). Para cada estado preparado, los núcleos
\(\mathcal I_A^a(A\mid x)\) y \(\mathcal I_B^b(B\mid x)\) expresan las
probabilidades publicadas por los lectores locales.

La forma local factorizada de la distribución conjunta es
\begin{equation}
 P_{\rm loc}(A,B\mid a,b)
 =\int_{X_{\rm enr}}
   \mathcal I_A^a(A\mid x)\,
   \mathcal I_B^b(B\mid x)\,\dd\mu(x).
 \label{eq:bell-hmt-factorizado}
\end{equation}
Esta ecuación conserva un estado global común y permite que dicho estado posea una
memoria arbitrariamente rica. Su afirmación característica es la factorización de
las respuestas una vez fijado \(x\). La independencia de los ajustes exige además
que la medida \(\mu\) sea común a todas las elecciones de \(a\) y \(b\), y el
muestreo completo exige que la población observada represente esa misma preparación
en todos los pares de ajustes.

Esta distinción se enlaza directamente con la definición HMT de medición. El
acoplamiento reversible entre sistema y aparato produce un estado conjunto; el
instrumento selecciona una partición de ese estado y el lector publica una de sus
fibras. Para un ajuste \(a\) y un resultado \(A\), la fibra
\[
 \Omega_{a,A}
 =\{x\in X_{\rm enr}:\operatorname{Out}_a(x)=A\}
\]
representa el conjunto de historias compatibles con la lectura. Ajustes diferentes
inducen particiones diferentes de un mismo espacio; la variable aleatoria publicada
queda, por tanto, tipada por su contexto instrumental. Así se formula la
contextualidad en el mismo lenguaje de instrumentos, fibras y supervivencia.

\subsection{Factorización local, cota CHSH y estabilidad bajo enriquecimiento}

Consideremos primero resultados binarios
\(\mathcal O_A=\mathcal O_B=\{-1,+1\}\) y dos ajustes por ala. Las
correlaciones y el funcional de Clauser--Horne--Shimony--Holt \cite{ClauserHorneShimonyHolt1969} son
\begin{equation}
 E_{ab}=\sum_{A,B=\pm1}AB\,P(A,B\mid a,b),
 \qquad
 S_{\rm CHSH}=E_{00}+E_{01}+E_{10}-E_{11}.
 \label{eq:chsh}
\end{equation}

\begin{theorem}[Cota CHSH bajo factorización local]
\label{thm:chsh-local}
Si la distribución conjunta posee la forma
\eqref{eq:bell-hmt-factorizado}, la preparación es independiente de los ajustes y
el muestreo preserva la población común en los cuatro pares de ajustes, entonces
\[
 \lvert S_{\rm CHSH}\rvert\leq2.
\]
\end{theorem}

\begin{proof}
Para cada \(x\), definamos las esperanzas condicionales
\[
 \overline A_a(x)=\sum_{A=\pm1}A\,\mathcal I_A^a(A\mid x),
 \qquad
 \overline B_b(x)=\sum_{B=\pm1}B\,\mathcal I_B^b(B\mid x).
\]
Ambas pertenecen a \([-1,1]\). El integrando de
\(S_{\rm CHSH}\) es
\[
 \overline A_0(\overline B_0+\overline B_1)
 +\overline A_1(\overline B_0-\overline B_1).
\]
Para cualesquiera \(u,v\in[-1,1]\), se cumple
\(\lvert u+v\rvert+\lvert u-v\rvert\leq2\). El valor absoluto del
integrando no excede, por tanto, \(2\). La integración respecto de una única
medida de probabilidad \(\mu\) conserva la cota.
\end{proof}

El teorema conduce a una consecuencia específicamente importante para HMT.

\begin{corollary}[Persistencia de la cota bajo enriquecimiento del estado]
\label{cor:no-go-bell-memoria}
La sustitución de una variable oculta elemental por el estado completo
\(x\in X_{\rm enr}\) conserva la cota
\(\lvert S_{\rm CHSH}\rvert\leq2\) siempre que los dos instrumentos sigan
factorizando como en \eqref{eq:bell-hmt-factorizado}.
\end{corollary}

\begin{proof}
El estado \(x\) desempeña el papel de la variable compartida en el
\cref{thm:chsh-local}. El argumento depende de la factorización de las respuestas
condicionales y de una preparación común a los cuatro pares de ajustes; la dimensión
interna de \(X_{\rm enr}\) deja intactas ambas propiedades.
\end{proof}

Este corolario fija positivamente el lugar donde actúa una teoría HMT de las
correlaciones cuánticas. La genealogía global proporciona el dominio común. Para
realizar una violación, el morfismo físico debe aportar una regla conjunta no
separable. En esta formulación, la correlación de Bell expresa una propiedad de la
representación conjunta del estado, compatible con la independencia operacional de
los marginales.

\subsection{Extensión de la entropía binaria al sistema qutrit y emergencia del invariante \(\log 3\)}

El TRIT selecciona de manera natural el alfabeto
\(\mathbb Z/3\mathbb Z\). Para una distribución
\(p=(p_0,p_1,p_2)\), su entropía expresada en trits es
\begin{equation}
 H_3(p)=-\sum_{j=0}^{2}p_j\log_3p_j
       =\frac{H_e(p)}{\log 3}.
 \label{eq:entropia-base-tres}
\end{equation}
La misma relación vale para la información mutua; para evitar confundirla con el
funcional de Bell que aparecerá enseguida, la denotaremos por
\[
 \mathsf I^{(3)}(X:Y)
 =\frac{\mathsf I^{(e)}(X:Y)}{\log 3}.
\]

\begin{proposition}[Invariancia de la localidad bajo cambio de unidad informativa]
\label{prop:log3-localidad}
El paso de logaritmos naturales a logaritmos en base tres multiplica entropías e
informaciones por la constante positiva \((\log 3)^{-1}\). Conserva las
distribuciones conjuntas, sus marginales y la condición de factorización local.
\end{proposition}

\begin{proof}
La identidad \(\log_3u=(\log u)/(\log 3)\) demuestra la primera afirmación.
La segunda se sigue porque el cambio de base actúa sobre el valor con el que se
cuantifica una distribución ya dada y deja invariantes tanto sus términos como los
núcleos estocásticos que la factorizan.
\end{proof}

La aparición de \(\log 3\) sitúa el análisis en una unidad informativa adecuada al
alfabeto ternario. El cambio sustantivo para Bell es el paso de dos a tres resultados
por instrumento. Ese paso conduce del funcional binario CHSH a una desigualdad
qutrit y permite que la estructura TRIT intervenga tanto en la definición de los
ajustes como en la elección de unidades.

\subsection{La desigualdad CGLMP para \(3\) resultados}

Sean \(A_1,A_2,B_1,B_2\) variables con valores en
\(\mathbb Z/3\mathbb Z\); todas las igualdades de esta sección se entienden
módulo tres. Definimos el funcional de
Collins--Gisin--Linden--Massar--Popescu \cite{CGLMP2002}
\begin{align}
 I_3^{\rm CGLMP}={}&P(A_1=B_1)+P(B_1=A_2+1)
 +P(A_2=B_2)+P(B_2=A_1)\notag\\
 &-P(A_1=B_1-1)-P(B_1=A_2)
 -P(A_2=B_2-1)-P(B_2=A_1-1).
 \label{eq:cglmp-tres}
\end{align}

\begin{theorem}[Cota local triádica]
\label{thm:cglmp-local}
Toda distribución de cuatro instrumentos ternarios que admita un modelo local
factorizado satisface
\[
 I_3^{\rm CGLMP}\leq2.
\]
\end{theorem}

\begin{proof}
El conjunto de distribuciones locales es la envolvente convexa de las asignaciones
deterministas, de modo que basta probar la cota en sus vértices. Para una asignación
fija, escribamos
\[
 x=A_1-B_1,\qquad y=B_1-A_2,\qquad
 z=A_2-B_2,\qquad t=B_2-A_1.
\]
La relación de cierre es \(x+y+z+t=0\). En estos términos,
\[
\begin{aligned}
 I_3^{\rm CGLMP}={}&
 \mathbf1_{x=0}+\mathbf1_{y=1}
 +\mathbf1_{z=0}+\mathbf1_{t=0}\\
 &-\mathbf1_{x=2}-\mathbf1_{y=0}
 -\mathbf1_{z=2}-\mathbf1_{t=2}.
\end{aligned}
\]
Las veintisiete elecciones de \((x,y,z)\) determinan \(t\). Para
\(y=0\), los nueve patrones producen una vez \(-4\), siete veces \(-1\) y
una vez \(2\); para \(y=1\), producen tres veces \(-1\) y seis veces \(2\);
para \(y=2\), producen seis veces \(-1\) y tres veces \(2\). El máximo en
todos los vértices es \(2\), y la convexidad conserva esa cota.
\end{proof}

La desigualdad anterior describe la frontera convexa de los modelos locales en el
mismo alfabeto que utiliza HMT. Su evaluación física parte de una distribución conjunta obtenida de una
preparación y de cuatro instrumentos concretos. La condición
\(I_3^{\rm CGLMP}>2\), cuando se obtiene con una muestra completa y ajustes fijados
independientemente, excluye una factorización local del tipo
\eqref{eq:bell-hmt-factorizado} para resultados ternarios.

\subsection{Transformada triádica desde la estructura elíptica interna}

La realización de los ajustes qutrit requiere una transformación que mezcle las
tres salidas. Sea \(E\) un espacio real euclídeo provisto de una estructura
compleja ortogonal \(J_E\), es decir,
\[
 J_E^2=-I,
 \qquad
 J_E^{\!*}=-J_E.
\]
Definamos
\[
 \omega_J=\exp\!\left(\frac{2\pi}{3}J_E\right),
 \qquad
 \omega_J^3=I,
 \qquad
 I+\omega_J+\omega_J^2=0.
\]
Sobre \(E\otimes\mathbb R^3\), la transformada triádica interna es
\begin{equation}
 \mathcal F_3^{(J)}
 =\frac1{\sqrt3}\sum_{r,s=0}^{2}
   \omega_J^{rs}\otimes|r\rangle\langle s|.
 \label{eq:fourier-triadica-interna}
\end{equation}

\begin{proposition}[Ortogonalidad de la transformada triádica interna]
\label{prop:fourier-triadica-unitaria}
La aplicación \(\mathcal F_3^{(J)}\) es ortogonal y su inversa se obtiene
cambiando \(J_E\) por \(-J_E\):
\[
 \bigl(\mathcal F_3^{(J)}\bigr)^{-1}
 =\mathcal F_3^{(-J)}.
\]
En una carta que identifica \(J_E\) con la multiplicación por \(i\), recupera la
matriz de Fourier compleja de orden tres.
\end{proposition}

\begin{proof}
La ortogonalidad de \(J_E\) implica
\((\omega_J^k)^{\!*}=\omega_J^{-k}\). El bloque \((r,r')\) de
\(\mathcal F_3^{(J)}\mathcal F_3^{(-J)}\) es
\[
 \frac13\sum_{s=0}^{2}\omega_J^{(r-r')s}.
\]
La suma vale \(I\) si \(r=r'\) y vale cero en los otros dos casos por
\(I+\omega_J+\omega_J^2=0\). La composición es, por tanto, la identidad.
La última afirmación resulta de
\(\exp(\theta J_E)\leftrightarrow e^{i\theta}\) en la carta compleja.
\end{proof}

Esta construcción hace operativo el tipo elíptico del TRIT en los instrumentos de
Bell. Un ajuste local puede componerse a partir de
\(\mathcal F_3^{(J)}\), transportes diagonales de fase y proyectores sobre la base
ternaria. Los transportes de orientación y memoria de cada ala han de determinar
esas fases antes de evaluar la desigualdad. De este modo, \(J_E\) aporta la rotación
interna, mientras el registro genealógico aporta la selección genealógica de los instrumentos.

\subsection{Representación conjunta e independencia de los marginales}

La confluencia con la mecánica cuántica se expresa mediante un mapa de realización
\begin{equation}
 \mathcal R_{\rm Bell}:X_{\rm enr}\longrightarrow
 \mathcal D(\mathcal H_A\otimes\mathcal H_B),
 \label{eq:realizacion-bell}
\end{equation}
donde \(\mathcal D(\mathcal H_A\otimes\mathcal H_B)\) es el espacio de
operadores densidad. Sean
\(\{M_{A,k}^{a}\}_{k\in\mathbb Z/3\mathbb Z}\) y
\(\{N_{B,\ell}^{b}\}_{\ell\in\mathbb Z/3\mathbb Z}\) POVM locales. La
regla conjunta es
\begin{equation}
 P(k,\ell\mid a,b)
 =\int_{X_{\rm enr}}
  \operatorname{tr}\!\left[
   \mathcal R_{\rm Bell}(x)
   \bigl(M_{A,k}^{a}\otimes N_{B,\ell}^{b}\bigr)
  \right]\dd\mu(x).
 \label{eq:born-bell-hmt}
\end{equation}

\begin{proposition}[Independencia de los marginales por completitud local]
\label{prop:no-signaling-hmt}
Si
\[
 \sum_{\ell=0}^{2}N_{B,\ell}^{b}=I_B
 \quad\text{para todo }b,
\]
entonces el marginal de Alice es independiente del ajuste de Bob. La afirmación
simétrica vale para la otra ala.
\end{proposition}

\begin{proof}
Al sumar \eqref{eq:born-bell-hmt} sobre \(\ell\), la completitud reemplaza
\(\sum_\ell N_{B,\ell}^{b}\) por \(I_B\). Se obtiene
\[
 \sum_{\ell}P(k,\ell\mid a,b)
 =\int\operatorname{tr}\!\left[
   \mathcal R_{\rm Bell}(x)
   \bigl(M_{A,k}^{a}\otimes I_B\bigr)
  \right]\dd\mu(x),
\]
expresión independiente de \(b\).
\end{proof}

La ecuación \eqref{eq:realizacion-bell} fija el tipo del entrelazador entre la
genealogía HMT y el espacio bipartito de la realización física. Cuando
\(\mathcal R_{\rm Bell}(x)\) es separable para casi todo \(x\), la distribución
puede reescribirse mediante variables compartidas y permanece en el dominio local.
Una realización no separable puede abandonar esa factorización y conservar, por la
proposición anterior, la independencia de los marginales respecto del ajuste remoto.
Por eso correlación global y transmisión de señal son propiedades diferentes de la
misma construcción.

\subsection{Reconstrucción holográfica de la correlación mediante clases de rutas indistinguibles}

La estructura HMT organiza ahora la cuestión local--global mediante objetos tipados.
La totalidad es una sección compatible del sistema de
extensiones, con memoria de su genealogía. Cada ala experimental es una frontera
dotada de un instrumento que publica una fibra. La preparación conjunta fija el
estado común; los instrumentos fijan las preguntas locales; la regla de Born fija
las probabilidades; CHSH o CGLMP determinan si esas probabilidades admiten una
descomposición local.

La correlación holográfica se realiza como lectura de un mismo estado global dentro
de una representación conjunta. La independencia de los marginales asegura que cada
estadística local permanece invariante frente al ajuste remoto, mientras la no
separabilidad puede
impedir que la distribución completa se reduzca a respuestas independientes
condicionadas por una variable clásica.

Queda así establecida una jerarquía probatoria precisa. Las cotas locales CHSH y
CGLMP, la persistencia de la cota bajo enriquecimiento del estado, la invariancia
bajo cambio a \(\log 3\), la transformada \(\mathcal F_3^{(J)}\) y la independencia
de los marginales por completitud de los POVM son resultados matemáticos cerrados.
La arquitectura HMT de
Bell está fijada por
\[
 X_{\rm enr}
 \xrightarrow{\ \mathcal R_{\rm Bell}\ }
 \mathcal D(\mathcal H_A\otimes\mathcal H_B)
 \xrightarrow{\ M^a\otimes N^b\ }
 P(k,\ell\mid a,b)
 \xrightarrow{\ \mathrm{CGLMP}\ }
 I_3^{\rm CGLMP}.
\]
La evaluación de una violación HMT concreta se obtiene al construir
\(\mathcal R_{\rm Bell}\) y los cuatro instrumentos desde el registro genealógico enriquecido,
certificar su independencia de ajustes y calcular la distribución sin
postselección. El capítulo proporciona el dominio, los operadores, las cotas y el
criterio de realización; la evaluación numérica queda reservada a la distribución
derivada de esos instrumentos.
